\documentclass[10pt,a4paper]{amsart}
\usepackage[margin=19mm]{geometry}
\usepackage{amsmath,amssymb,mathtools}
\usepackage[scr=rsfs,cal=euler]{mathalfa}
\usepackage{xfrac}
\usepackage[unicode,hidelinks,bookmarks=false]{hyperref}
\newcommand{\ie}{i.e.\ }
\newcommand{\cf}{cf.\ }
\newcommand{\resp}{respectively\ }
\newcommand{\Subsection}{\subsection}
\providecommand{\nobreakdash}{\nobreak\hbox{-}}
\setlength{\emergencystretch}{2em}
\multlinegap0pt
\usepackage{enumerate}
\usepackage{xcolor}
\newtheorem{lemma}{Lemma}
\title{On the Ban-Linial Conjecture}
\author{Matt DeVos, Kathryn Nurse}
\date{Source: arXiv:2512.18913v1 (CC BY 4.0)}
\begin{document}
\maketitle

\noindent\emph{Source and licence.} On the Ban-Linial Conjecture. Version arXiv:2512.18913v1 (CC BY 4.0), distributed by arXiv under the Creative Commons Attribution 4.0 International licence, http://creativecommons.org/licenses/by/4.0/, as recorded in the arXiv licence metadata for this exact version and shown on the abstract page https://arxiv.org/abs/2512.18913v1. Full licence notice: \texttt{licenses/arxiv-2512.18913-CC-BY-4.0.txt}.\par\smallskip

\noindent\textit{Editorial scope.} Counted unit: the article's lemma unrootedlemma (source0.tex lines 147--155). The article's complete original proof of it is delivered with it (source0.tex lines 157--188, one \texttt{proof} environment of 32 lines). No mathematics is written, repaired or re-derived: the statement and the proof are the article's own lines, in the article's own notation, and no retained line is substituted or re-flowed.  No further source-local context is needed: every cross-reference of every retained line resolves inside the two delivered spans.  One declared adaptation: exactly 2 ancillary strings inside the retained spans are omitted (image-inclusion commands and, where the source repeats a label, the repeated label definitions) and no other line differs from the source; the figure environments, with their placement, captions and labels, are kept, and the package records the omitted strings in fidelity receipts bound to the affected bodies.  No retained line contains a citation, so the wrapper carries no bibliography and no bibliography entry.  The retained lines contain the article's own diagram code, which the wrapper typesets with the standard tikz packages; no image file is delivered and the package declares no asset dependency.  Editorial formatting only, in addition: an AMS \texttt{amsart} preamble that restates the article's own declarations and macros used by the retained lines, namely the environments \texttt{lemma} on separate counters, together with the standard packages those lines need.  The retained environments print in this document as Lemma 1 in order of appearance, whereas the article prints the same items with the numbering of its own full document.  The complete native source of the article as distributed in the arXiv e-print for this exact version is bundled unchanged as \texttt{sources/191432/source0.tex}.
\begin{lemma}
\label{unrootedlemma}
Let $T$ be a cubic tree with $|V(T)| \ge 4$, let $(X,Y)$ be a split of $V_1(G)$, and let $\epsilon = \pm 1$.  There exists a split $(X^*,Y^*)$ of $T$ satisfying all of the following properties:
\begin{enumerate}
\item $X \subseteq X^*$ and $Y \subseteq Y^*$,
\item $\mathrm{disc}(X^*,Y^*) \in \epsilon \{0,1,2\}$
\item $T - E(X^*,Y^*)$ has at most one vertex of degree greater than one.
\end{enumerate}
\end{lemma}

\begin{proof}
We proceed by induction on $|V|$ with the base case $|V| =4$ holding by inspection.  Next suppose there exists a pair of cherries $\{w,w'\}$ both adjacent to the vertex $v \in V$ so that $w,w' \in X$.  In this case the result follows by applying induction to the tree $T - \{w,w'\}$ and the sets $X' = X \setminus \{w,w'\}$ and $Y' = Y \cup \{v\}$.  A similar argument holds if $w,w' \in Y$, so we may assume that we do not have a pair of cherries of the same colour.  If $|V| = 6$ or $|V| = 8$, then up to isomorphism and colour swapping there is just one possibility for $T$, and in both of these cases the result is straightforward to verify.  So we may now assume $|V| \ge 10$.  Next we will utilize some reductions that we describe with figures.  We will treat our splits $(X,Y)$ and $(X^*,Y^*)$ as (improper) colourings of the vertices, and we adopt the convention for our figures that a solid black vertex is in $X$ and an open white vertex is in $Y$.  

\bigskip

\noindent{\it Reduction 1.}

If $T$ contains a vertex $v$ with two neighbours both of which are adjacent with a pair of cherries, then we perform reduction 1 to reduce to a smaller tree $T'$ and flip the sign of $\epsilon$ as shown in the figure below.

\begin{center}

\end{center}

Now we apply induction to colour the vertices of the smaller tree.  If we have the case on the left (right) in the above figure, we extend back to a colouring of $T$ by assigning both neighbours of $v$ in $V(T) \setminus V(T')$ the colour black (white).  This operation either increases or decreases the discrepancy by 2, but in either case we have the desired colouring of $T$.  


\bigskip

\noindent{\it Reduction 2.}

If $T$ contains a vertex $v$ adjacent to a leaf vertex and another vertex $w$ so that $w$ is adjacent with a pair of cherries, then we perform reduction 2 to reduce $T$ to a smaller tree $T'$

\begin{center}

\end{center}

Now we apply induction to colour the vertices of this smaller tree $T'$ using the original parameter $\epsilon$.  Let $u$ be the unique neighbour of $v$ that is not a leaf and is not equal to $w$ (as shown in the figure).  Since the two versions of reduction 2 are similar, we will assume that we have applied the one on the left (so $v$ is coloured white in $T'$).  If $u$ is coloured black in the resulting colouring of $T'$ then we keep this colour for $u$ and assign $w$ the colour black.  The resulting colouring of $T$ has the same discrepancy as that of $T'$ and will thus satisfy the lemma.  If $u$ is coloured white in the resulting colouring of $T'$ then we recolour $v$ to be black and assign $w$ the colour white.  Again the resulting colouring of $T$ has the same discrepancy as that of $T'$ and this satisfies the lemma.  

\bigskip

To complete the proof, it now suffices to show that one of our two reductions can be performed.  To see this, recall that every pair of cherries in $T$ have opposite colours, one black and one white.  As a bookkeeping device, for every pair of cherries with one white and the other black, let $v$ be the common neighbour, delete these cherries and assign $v$ the colour red.   After this operation has been performed wherever possible, the resulting tree is still cubic, and must have a pair of cherries at least one of which is coloured red.  In all cases, this pair of cherries gives us a reduction of one of our two types.
\end{proof}

\end{document}
